\subsection{Работа с модулями платформы}
\begin{enumerate}
    \item \verb|Ledbar.new(N)| — создание объекта управления (обычно \verb|N=25|).
    \item \verb|leds:set(i, r, g, b)| — установка цвета светодиода \verb|i|. Параметры \verb|r,g,b| — от 0.0 до 1.0.
    \item \verb|Timer.new(period, fn)| — создание периодического таймера.
    \item \verb|Timer.callLater(delay, fn)| — отложенный вызов функции.
    \item \verb|function callback(event)| — обработчик системных событий.
\end{enumerate}

\subsection{Типичные ошибки новичка}
\begin{enumerate}
    \item \textbf{Регистр букв}: \texttt{Print}, \texttt{PRINT} и \texttt{print} — это три разных слова. В Lua все ключевые слова и функции пишутся строчными буквами (кроме классов вроде \texttt{Ledbar}).
    \item \textbf{Путаница с индексами}:
    \begin{itemize}
        \item Обычные таблицы Lua нумеруются с \textbf{1}: \texttt{t[1]}.
        \item Аппаратные модули (светодиоды, массивы из C++) часто нумеруются с \textbf{0}: \texttt{leds:set(0, ...)}. Будьте внимательны!
    \end{itemize}
    \item \textbf{Бесконечные циклы}:
    \begin{codefile}{lua/examples/chapter_04/lua_platform_infinite_loop.lua}
    \end{codefile}
    Вместо вечных циклов используйте \texttt{Timer}.
    \item \textbf{Ошибка "attempt to index a nil value"}: Вы пытаетесь обратиться к полю переменной, которая пуста (\texttt{nil}). Проверьте, правильно ли вы создали объект или таблицу.
\end{enumerate}

\subsection{Примеры, используемые далее}
\begin{infobox}[Обратите внимание]
    В следующих примерах используется библиотека \texttt{Ledbar}, которую мы подробно изучим в следующей главе. Сейчас сосредоточьтесь на синтаксисе языка Lua (циклы, функции, таблицы), не вдаваясь в детали работы светодиодов.
\end{infobox}
\begin{codefile}{lua/examples/chapter_04/lua_platform_leds_unpack.lua}
\end{codefile}
\begin{codefile}{lua/examples/chapter_04/lua_platform_timer.lua}
\end{codefile}
\begin{codefile}{lua/examples/chapter_04/lua_platform_blink.lua}
\end{codefile}
